\documentclass{article}
\usepackage[utf8]{inputenc}
\usepackage{amsmath}
\usepackage{geometry}
\geometry{margin=2cm}
\begin{document}

\section*{Análise da Questão 163 — ENEM 2024 — Matemática}

\subsection*{Enunciado}
Um sistema de polias circulares e correias é um dos mecanismos responsáveis pela transmissão de movimento em máquinas rotativas. O manual de um motor apresenta um sistema com duas polias e uma correia tensionada e ajustada para não apresentar folgas, com segmentos de reta tangentes às polias onde não há contato direto.

Para substituir essa correia, é necessário especificar seu comprimento, sabendo que $\pi$ será aproximado por 3.

\subsection*{Solução e Análise}

\input{\detokenize{\text{(HERE GOES THE CONTENT OF \textbackslash technicalSolution)}}}

\subsection*{Fórmulas Utilizadas}

\input{\detokenize{\text{(HERE GOES THE CONTENT OF \textbackslash formulasUsed)}}}

\subsection*{Evidências Visuais}

\input{\detokenize{\text{(HERE GOES THE CONTENT OF \textbackslash visualEvidenceUsed)}}}

\subsection*{Explicação Didática}

\input{\detokenize{\text{(HERE GOES THE CONTENT OF \textbackslash didacticExplanation)}}}

\subsection*{Erros Comuns}

\input{\detokenize{\text{(HERE GOES THE CONTENT OF \textbackslash commonPitfalls)}}}

\subsection*{Dicas para Ensino}

\input{\detokenize{\text{(HERE GOES THE CONTENT OF \textbackslash teachingTips)}}}

\subsection*{Distratores}

\input{\detokenize{\text{(HERE GOES THE CONTENT OF \textbackslash distractors)}}}

\subsection*{Perfil da Questão}

\input{\detokenize{\text{(HERE GOES THE CONTENT OF \textbackslash itemProfile)}}}

\noindent \textit{Observação: A imagem do enunciado apresenta as duas polias de raios $4$ cm e $8$ cm, distância entre centros de $15$ cm e o ângulo de contato de $150^\circ$ na polia menor, conforme descrito nos dados acima, com os segmentos de correia tangentes representados por linhas retas entre os pontos de contato nas polias.}

% Para compilar o documento completo, substituir os comandos \input com o conteúdo respectivo fornecido nos campos do JSON.

\end{document}